\documentclass[11pt,a4paper]{article}

% ---------- Codificación y tipografía ----------
\usepackage[utf8]{inputenc}
\usepackage[T1]{fontenc}
\usepackage{lmodern}
\usepackage[spanish,es-nodecimaldot]{babel}
\usepackage{microtype}
\usepackage{amsmath,amssymb}
\usepackage{geometry}
\geometry{top=2.3cm,bottom=2.3cm,left=2.2cm,right=2.2cm}
\usepackage{enumitem}
\usepackage{booktabs}
\usepackage{longtable}
\usepackage{array}
\usepackage{fancyhdr}
\usepackage{titlesec}
\usepackage{setspace}
\usepackage{multicol}
\usepackage{needspace}

\setstretch{1.08}
\setlength{\parindent}{0pt}
\setlength{\parskip}{4pt}

% ---------- Encabezado y pie ----------
\pagestyle{fancy}
\fancyhf{}
\lhead{\small Banco de Preguntas PEU --- SUNASS}
\rhead{\small Módulo III: Competencia Perfecta}
\cfoot{\small \thepage}
\renewcommand{\headrulewidth}{0.4pt}
\renewcommand{\footrulewidth}{0pt}

% ---------- Títulos ----------
\titleformat{\section}{\large\bfseries}{\thesection}{0.8em}{}
\titlespacing*{\section}{0pt}{14pt}{6pt}
\titleformat{\subsection}{\normalsize\bfseries}{\thesubsection}{0.6em}{}
\titlespacing*{\subsection}{0pt}{10pt}{4pt}

% ---------- Comandos del examen ----------
\newcounter{qnum}
\newcommand{\q}[1]{%
  \needspace{8\baselineskip}%
  \par\medskip\refstepcounter{qnum}%
  \noindent\textbf{\theqnum.}\ #1\par\nopagebreak[4]
}
\newenvironment{opts}
  {\begin{enumerate}[label=\Alph*),leftmargin=2.0em,itemsep=1pt,topsep=2pt,parsep=0pt,partopsep=0pt]}
  {\end{enumerate}}
\makeatletter
\@itempenalty=10000
\makeatother
\newcommand{\sol}[3]{\par\noindent\textbf{#1. #2.}\ #3\par\smallskip}

% ==============================================================================
% Configuración BackBone Academy para los Bancos de Preguntas (Módulos 01-10)
% ==============================================================================
\usepackage{graphicx}
\usepackage{xcolor}
\usepackage{tikz}
\usepackage[hidelinks]{hyperref}

% ==============================================================================
% Identidad de marca BackBone Academy compartida (Beamer y Bancos de Preguntas)
% ==============================================================================
\definecolor{azulPrincipal}{RGB}{43, 45, 162}
\definecolor{azulOscuro}{RGB}{24, 26, 98}
\definecolor{grisTexto}{RGB}{85, 85, 95}

\newcommand{\backboneIgUrl}{https://www.instagram.com/backbone.academy}
\newcommand{\backboneTtUrl}{https://www.tiktok.com/@backbone.academy}
\newcommand{\backboneWebUrl}{https://leonardovelasquez14-detg.github.io/peu-sunass-economia/}

% Icono discreto de Instagram en TikZ
\newcommand{\igicon}[1][white!88]{%
  \tikz[baseline=-0.25ex]{%
    \draw[#1, line width=0.65pt, rounded corners=1.1pt] (0,0) rectangle (0.23cm,0.23cm);
    \draw[#1, line width=0.6pt] (0.115cm,0.115cm) circle (0.055cm);
    \fill[#1] (0.18cm,0.18cm) circle (0.014cm);
  }%
}

% Icono discreto de TikTok en TikZ (nota musical estilizada)
\newcommand{\tiktokicon}[1][white!88]{%
  \tikz[baseline=-0.25ex]{%
    \draw[#1, line width=0.82pt, line cap=round, line join=round]
      (0.11cm,0.06cm) -- (0.11cm,0.22cm)
      .. controls (0.12cm,0.16cm) and (0.16cm,0.145cm) .. (0.20cm,0.145cm);
    \draw[#1, line width=0.82pt, line cap=round]
      (0.11cm,0.065cm) arc[start angle=0, end angle=295, radius=0.048cm];
  }%
}

% Enlaces compactos de redes sociales: \backboneSocials{color}{tamaño}
\newcommand{\backboneSocials}[2]{%
  #2\color{#1}%
  \href{\backboneIgUrl}{\igicon[#1]}%
  \hspace{3.5pt}%
  \href{\backboneTtUrl}{\tiktokicon[#1]}%
  \hspace{4pt}%
  \href{\backboneIgUrl}{\textbf{@backbone.academy}}%
}


% Pie de página institucional en todas las páginas interiores
\lfoot{\scriptsize\color{azulOscuro}\textbf{BackBone Academy}\ \color{grisTexto}· \href{\backboneIgUrl}{\igicon[azulPrincipal]\ \textbf{@backbone.academy}} · \href{\backboneTtUrl}{\tiktokicon[azulPrincipal]\ \textbf{TikTok}} · \href{\backboneWebUrl}{Simulador Web}}
\cfoot{}
\rfoot{\small\thepage}
\renewcommand{\footrulewidth}{0.3pt}

% Encabezado superior para la portada de cada banco
\newcommand{\backboneportadaheader}{%
  \vspace*{0.1cm}
  \noindent\includegraphics[height=0.72cm]{../teoria/backbone_logo_azul.png}%
  \hfill%
  \raisebox{0.22cm}{\backboneSocials{azulPrincipal}{\small}}%
  \par\vspace{0.9cm}%
}

% Pie institucional para la portada de cada banco
\newcommand{\backboneportadafooter}{%
  {\small\color{azulOscuro}\textbf{BackBone Academy} \color{black}· Material de preparación académica.}\\[0.18cm]
  {\small\color{azulPrincipal}\href{\backboneIgUrl}{\igicon[azulPrincipal]\ \textbf{Instagram}} \quad·\quad \href{\backboneTtUrl}{\tiktokicon[azulPrincipal]\ \textbf{TikTok}} \textbf{(@backbone.academy)} \quad·\quad \href{\backboneWebUrl}{\textbf{Simulador Web Interactivo}}}\\[0.18cm]
  {\small Edición 2026}%
}


\begin{document}

% ============================================================
% PORTADA
% ============================================================
\begin{titlepage}
\thispagestyle{empty}
\centering
\backboneportadaheader

{\Large\bfseries BANCO DE PREGUNTAS}\\[0.35cm]
{\large EXAMEN DE ADMISIÓN}\\[0.15cm]
{\large PROGRAMA DE EXTENSIÓN UNIVERSITARIA (PEU)}\\[0.15cm]
{\large SUPERINTENDENCIA NACIONAL DE SERVICIOS DE SANEAMIENTO}\\[0.05cm]
{\large SUNASS}\\[1.0cm]

\rule{\linewidth}{0.9pt}\\[0.55cm]
{\LARGE\bfseries MÓDULO III: COMPETENCIA PERFECTA}\\[0.35cm]
{\large Especialidad: Economía}\\[0.55cm]
\rule{\linewidth}{0.9pt}\\[1.0cm]

\begin{tabular}{ll}
\toprule
Número de preguntas: & 90 \\
Bloque A: & 30 preguntas \\
Bloque B: & 35 preguntas \\
Bloque C: & 25 preguntas \\
Duración sugerida: & 150 minutos \\
Nivel: & Microeconomía intermedia--avanzada \\
Formato: & Opción múltiple (A, B, C, D) \\
\bottomrule
\end{tabular}

\vfill
\backboneportadafooter
\end{titlepage}

% ============================================================
% INSTRUCCIONES
% ============================================================
\section*{Instrucciones generales}

\begin{enumerate}[label=\arabic*.,leftmargin=1.6em,itemsep=3pt]
  \item El presente banco contiene \textbf{90 preguntas de opción múltiple}. Cada pregunta tiene una sola respuesta correcta.
  \item No se descuenta por respuesta incorrecta; se recomienda responder todas las preguntas.
  \item El tiempo sugerido de resolución es de \textbf{150 minutos}. Use calculadora no programable.
  \item Las preguntas están organizadas en tres bloques: fundamentos teóricos, aplicaciones cuantitativas y casos regulatorios del sector saneamiento.
  \item El \textbf{solucionario completo} y la \textbf{tabla de claves rápidas} se encuentran al final. Se recomienda resolver el examen antes de consultarlos.
\end{enumerate}

\section*{Estructura del banco}

\begin{center}
\renewcommand{\arraystretch}{1.25}
\begin{tabular}{@{}p{1.4cm}>{\raggedright\arraybackslash}p{7.6cm}cc@{}}
\toprule
\textbf{Bloque} & \textbf{Contenido} & \textbf{Preguntas} & \textbf{Peso} \\
\midrule
A & Fundamentos teóricos y conceptuales: supuestos de competencia perfecta, maximización del beneficio, oferta individual y de la industria, equilibrio de corto y largo plazo, excedentes y eficiencia. & 1--30 & 33.3\% \\
B & Aplicaciones cuantitativas y optimización: funciones de costo y beneficio, oferta de la empresa y de la industria, equilibrio competitivo, entrada y salida, impuestos, subsidios y pérdida de bienestar. & 31--65 & 38.9\% \\
C & Casos con implicancias tarifarias y regulatorias del sector saneamiento: competencia y monopolio natural, precios máximos, incidencia de impuestos, subsidios, derechos de agua y excedentes. & 66--90 & 27.8\% \\
\bottomrule
\end{tabular}
\end{center}

\section*{Competencias evaluadas}

\noindent El banco evalúa la capacidad de interpretar los supuestos de la competencia perfecta, resolver problemas de maximización de beneficios, derivar ofertas individuales y de industria, determinar equilibrios de corto y largo plazo, calcular excedentes y evaluar los efectos de impuestos, subsidios y controles de precios en mercados de agua y saneamiento.

\section*{Notación principal}

{\small
\renewcommand{\arraystretch}{1.0}
\begin{center}
\begin{tabular}{@{}ll@{\hspace{1.0cm}}ll@{}}
\toprule
$Q$ & Cantidad de mercado & $q$ & Cantidad de la empresa \\
$P$ & Precio & $IT$ & Ingreso total \\
$CT$ & Costo total & $IMg$ & Ingreso marginal \\
$CMg$ & Costo marginal & $CMe$ & Costo medio \\
$CV$ & Costo variable & $CVMe$ & Costo variable medio \\
$CF$ & Costo fijo & $\pi$ & Beneficio económico \\
$EC$ & Excedente del consumidor & $EP$ & Excedente del productor \\
$ET$ & Excedente total & $PCP$ & Pérdida de bienestar \\
\bottomrule
\end{tabular}
\end{center}
}

\newpage
\setcounter{qnum}{0}

\section*{Bloque A. Fundamentos teóricos y conceptuales}

\noindent\textit{Preguntas 1 a 30. Evalúan la comprensión analítica de los supuestos de la competencia perfecta, la maximización del beneficio, la oferta individual y de la industria, el equilibrio de corto y largo plazo, los excedentes y la eficiencia asignativa.}

\q{¿Cuál de las siguientes combinaciones describe mejor un mercado de competencia perfecta?}
\begin{opts}
  \item Pocas empresas, producto diferenciado y barreras a la entrada.
  \item Una sola empresa, producto único y poder de mercado.
  \item Muchos compradores y vendedores, producto homogéneo, información perfecta y libre entrada y salida.
  \item El Estado fija el precio y las cantidades de todas las empresas.
\end{opts}

\q{La curva de demanda que enfrenta una empresa perfectamente competitiva es:}
\begin{opts}
  \item perfectamente elástica al precio de mercado.
  \item descendente y con pendiente negativa.
  \item vertical.
  \item quebrada.
\end{opts}

\q{La curva de demanda del mercado en competencia perfecta es:}
\begin{opts}
  \item horizontal.
  \item vertical.
  \item descendente.
  \item perfectamente elástica.
\end{opts}

\q{La condición de maximización del beneficio de una empresa es:}
\begin{opts}
  \item $P=CMe$.
  \item $IMg=CMe$.
  \item $P=CVMe$.
  \item $IMg=CMg$.
\end{opts}

\q{Para una empresa perfectamente competitiva, el ingreso marginal es:}
\begin{opts}
  \item mayor que el precio.
  \item menor que el precio.
  \item igual al precio.
  \item igual a cero.
\end{opts}

\q{El beneficio económico de una empresa se define como:}
\begin{opts}
  \item el ingreso total menos los costos explícitos.
  \item el beneficio contable.
  \item el ingreso total menos los costos fijos.
  \item el ingreso total menos el costo económico total, que incluye costos explícitos e implícitos.
\end{opts}

\q{El beneficio normal se relaciona con el beneficio económico porque:}
\begin{opts}
  \item el beneficio económico siempre es positivo.
  \item el beneficio normal es el costo de oportunidad del capital; un beneficio económico nulo significa que se está obteniendo el beneficio normal.
  \item el beneficio normal coincide con el beneficio contable.
  \item el beneficio normal es un costo fijo irrelevante.
\end{opts}

\q{En el corto plazo, una empresa competitiva puede obtener:}
\begin{opts}
  \item siempre beneficios positivos.
  \item siempre beneficios nulos.
  \item siempre pérdidas.
  \item beneficios positivos, nulos o negativos, según la relación entre el precio y el costo medio.
\end{opts}

\q{En el corto plazo, la empresa competitiva suspende la producción cuando:}
\begin{opts}
  \item $P<CMe$.
  \item $P<CVMe$.
  \item $P<CMg$.
  \item $P<CFMe$.
\end{opts}

\q{La curva de oferta de una empresa competitiva en el corto plazo es:}
\begin{opts}
  \item la curva de costo medio.
  \item la curva de costo variable medio.
  \item la curva de demanda.
  \item el tramo de la curva de costo marginal situado por encima del mínimo del costo variable medio.
\end{opts}

\q{La curva de oferta de la industria competitiva se obtiene mediante:}
\begin{opts}
  \item la suma vertical de las curvas de oferta individuales.
  \item la suma horizontal de las curvas de oferta individuales.
  \item el promedio de las curvas de oferta individuales.
  \item el máximo de las curvas de oferta individuales.
\end{opts}

\q{En el equilibrio de largo plazo de una industria perfectamente competitiva con libre entrada y salida:}
\begin{opts}
  \item $P>CMe$ y hay beneficios extraordinarios.
  \item $P<CVMe$.
  \item $P=CMg=CMe$ y el beneficio económico es cero.
  \item $P=CMg>CMe$.
\end{opts}

\q{Si en una industria competitiva el precio es mayor que el costo medio mínimo, en el largo plazo:}
\begin{opts}
  \item las empresas salen y el precio sube.
  \item no hay entrada de empresas.
  \item el precio sube por falta de oferta.
  \item entran empresas, la oferta aumenta y el precio tiende a bajar.
\end{opts}

\q{Si en una industria competitiva el precio es menor que el costo medio mínimo, en el largo plazo:}
\begin{opts}
  \item salen empresas, la oferta disminuye y el precio tiende a subir.
  \item entran empresas y la oferta aumenta.
  \item el precio sigue bajando indefinidamente.
  \item no cambia nada.
\end{opts}

\q{En una industria de costo constante, la curva de oferta de largo plazo es:}
\begin{opts}
  \item ascendente.
  \item descendente.
  \item horizontal.
  \item vertical.
\end{opts}

\q{En una industria de costo creciente, la curva de oferta de largo plazo es:}
\begin{opts}
  \item ascendente.
  \item horizontal.
  \item descendente.
  \item vertical.
\end{opts}

\q{En una industria de costo decreciente, la curva de oferta de largo plazo es:}
\begin{opts}
  \item horizontal.
  \item ascendente.
  \item vertical.
  \item descendente.
\end{opts}

\q{El excedente del productor se define como:}
\begin{opts}
  \item el ingreso total menos el costo total.
  \item el beneficio contable.
  \item el ingreso total menos el costo variable.
  \item el excedente del consumidor.
\end{opts}

\q{El excedente del consumidor se define como:}
\begin{opts}
  \item la disposición a pagar menos el precio efectivamente pagado.
  \item el precio menos el costo marginal.
  \item el ingreso total menos el costo total.
  \item la recaudación impositiva.
\end{opts}

\q{El excedente total o bienestar social es:}
\begin{opts}
  \item la suma del excedente del consumidor y del productor, y se maximiza en el equilibrio competitivo.
  \item únicamente la recaudación del impuesto.
  \item la diferencia entre el excedente del productor y el del consumidor.
  \item cero en el equilibrio de mercado.
\end{opts}

\q{Un precio máximo fijado por debajo del precio de equilibrio genera:}
\begin{opts}
  \item un excedente de oferta.
  \item ningún efecto sobre el mercado.
  \item escasez y pérdida de bienestar.
  \item un aumento automático de la oferta.
\end{opts}

\q{Un precio mínimo fijado por encima del precio de equilibrio genera:}
\begin{opts}
  \item un excedente de oferta y pérdida de bienestar.
  \item escasez.
  \item ningún efecto.
  \item un aumento de la demanda.
\end{opts}

\q{La incidencia económica de un impuesto depende de:}
\begin{opts}
  \item las elasticidades relativas de oferta y demanda; el lado menos elástico soporta una mayor carga.
  \item exclusivamente de quién paga legalmente el impuesto.
  \item siempre de los consumidores.
  \item siempre de los productores.
\end{opts}

\q{En el caso de un subsidio, la incidencia económica implica que:}
\begin{opts}
  \item siempre se beneficia el consumidor.
  \item el lado menos elástico captura una mayor parte del subsidio.
  \item siempre se beneficia el productor.
  \item no hay efectos sobre el precio.
\end{opts}

\q{La eficiencia asignativa en competencia perfecta se alcanza cuando:}
\begin{opts}
  \item $P=CMe$.
  \item $P=CVMe$.
  \item $P=CMg$.
  \item $P=0$.
\end{opts}

\q{La eficiencia productiva en el largo plazo competitivo se alcanza cuando:}
\begin{opts}
  \item $P>CMe$.
  \item $P=CMg>CMe$.
  \item $P=\min CMe$.
  \item $P=CVMe$.
\end{opts}

\q{La libre entrada de empresas en una industria competitiva:}
\begin{opts}
  \item aumenta los beneficios extraordinarios de las empresas existentes.
  \item no tiene efecto sobre los beneficios.
  \item crea un monopolio.
  \item elimina los beneficios económicos en el largo plazo.
\end{opts}

\q{Una empresa es precio-aceptante cuando:}
\begin{opts}
  \item fija un precio superior al de mercado.
  \item enfrenta una demanda descendente.
  \item puede influir en el precio mediante su nivel de producción.
  \item no puede influir en el precio de mercado y lo toma como dado.
\end{opts}

\q{El equilibrio de mercado en competencia perfecta se determina cuando:}
\begin{opts}
  \item $Q_d=Q_s$.
  \item $Q_d>Q_s$.
  \item $Q_d<Q_s$.
  \item $P=0$.
\end{opts}

\q{En el corto plazo, una empresa competitiva produce:}
\begin{opts}
  \item donde $P=CMg$, siempre que $P\geq CVMe$.
  \item siempre en el punto de cierre.
  \item donde $P=CMe$.
  \item fijando un precio superior al de mercado.
\end{opts}

\section*{Bloque B. Aplicaciones cuantitativas y optimización}

\noindent\textit{Preguntas 31 a 65. Requieren resolver problemas de maximización de beneficios, derivar ofertas individuales y de industria, calcular equilibrios de corto y largo plazo, determinar excedentes y evaluar los efectos de impuestos, subsidios y controles de precios.}

\q{Si $C(q)=100+10q+q^2$ y $P=30$, la cantidad que maximiza el beneficio es:}
\begin{opts}
  \item 5
  \item 10
  \item 15
  \item 20
\end{opts}

\q{En el problema anterior, el beneficio económico en el óptimo es:}
\begin{opts}
  \item 0
  \item 100
  \item $-100$
  \item 50
\end{opts}

\q{Si $C(q)=50+5q+0.5q^2$ y $P=25$, la cantidad que maximiza el beneficio es:}
\begin{opts}
  \item 10
  \item 15
  \item 20
  \item 25
\end{opts}

\q{En el problema anterior, el beneficio económico es:}
\begin{opts}
  \item 100
  \item 150
  \item 200
  \item 250
\end{opts}

\q{Si $C(q)=100+4q+q^2$ y $P=20$, el beneficio económico en el óptimo es:}
\begin{opts}
  \item $-36$
  \item 36
  \item 0
  \item $-64$
\end{opts}

\q{Con $C(q)=100+4q+q^2$, el precio de cierre de la empresa en el corto plazo es:}
\begin{opts}
  \item 0
  \item 2
  \item 24
  \item 4
\end{opts}

\q{Si $C(q)=q^2+10q+50$ y $P=30$, la cantidad ofrecida por la empresa es:}
\begin{opts}
  \item 5
  \item 10
  \item 15
  \item 20
\end{opts}

\q{Si 100 empresas idénticas tienen oferta individual $q=(P-10)/2$, la oferta de la industria es:}
\begin{opts}
  \item $Q=50P-500$
  \item $Q=50P+500$
  \item $Q=100P-1000$
  \item $Q=100P-500$
\end{opts}

\q{Si la demanda de mercado es $Q=1000-50P$ y la oferta es $Q=50P-500$, el precio de equilibrio es:}
\begin{opts}
  \item 15
  \item 10
  \item 20
  \item 25
\end{opts}

\q{En el equilibrio anterior, la cantidad de mercado es:}
\begin{opts}
  \item 100
  \item 150
  \item 200
  \item 250
\end{opts}

\q{En el mismo equilibrio, la cantidad producida por cada una de las 100 empresas es:}
\begin{opts}
  \item 1
  \item 2.5
  \item 5
  \item 10
\end{opts}

\q{Si $C(q)=q^2+4q+100$, el precio de equilibrio de largo plazo es:}
\begin{opts}
  \item 20
  \item 22
  \item 24
  \item 26
\end{opts}

\q{En el equilibrio de largo plazo anterior, la cantidad producida por cada empresa es:}
\begin{opts}
  \item 10
  \item 5
  \item 20
  \item 24
\end{opts}

\q{Si la demanda de mercado es $Q=500-10P$ y cada empresa produce 10 unidades en el largo plazo, el número de empresas es:}
\begin{opts}
  \item 20
  \item 24
  \item 26
  \item 30
\end{opts}

\q{En el equilibrio de largo plazo con $P=24$ y $Q=260$, el excedente del consumidor es:}
\begin{opts}
  \item 2600
  \item 3380
  \item 5000
  \item 6760
\end{opts}

\q{En el equilibrio de largo plazo competitivo con libre entrada, el beneficio económico de cada empresa es:}
\begin{opts}
  \item 3380
  \item 2600
  \item 0
  \item 260
\end{opts}

\q{Si en una industria de costo constante la demanda aumenta, en el nuevo equilibrio de largo plazo el precio:}
\begin{opts}
  \item vuelve a 24.
  \item sube permanentemente a 30.
  \item baja a 20.
  \item sube a 26 y se queda allí.
\end{opts}

\q{Si la demanda es $Q_d=120-2P$ y la oferta es $Q_s=3P-30$, el precio de equilibrio es:}
\begin{opts}
  \item 20
  \item 30
  \item 40
  \item 50
\end{opts}

\q{En el equilibrio anterior, la cantidad de equilibrio es:}
\begin{opts}
  \item 30
  \item 50
  \item 60
  \item 80
\end{opts}

\q{En el mismo mercado, el excedente del consumidor es:}
\begin{opts}
  \item 300
  \item 600
  \item 1200
  \item 900
\end{opts}

\q{En el mismo mercado, el excedente del productor es:}
\begin{opts}
  \item 300
  \item 900
  \item 600
  \item 1200
\end{opts}

\q{En el mismo mercado, el excedente total es:}
\begin{opts}
  \item 900
  \item 1200
  \item 600
  \item 1500
\end{opts}

\q{Si se fija un precio máximo de 20 en ese mercado, la escasez es:}
\begin{opts}
  \item 20
  \item 30
  \item 40
  \item 50
\end{opts}

\q{Con el precio máximo de 20, la pérdida de bienestar es:}
\begin{opts}
  \item 625
  \item 500
  \item 750
  \item 1250
\end{opts}

\q{Si se aplica un impuesto específico de 5 por unidad, el nuevo precio de equilibrio es:}
\begin{opts}
  \item 30
  \item 31
  \item 32
  \item 33
\end{opts}

\q{Con el impuesto de 5, la cantidad de equilibrio es:}
\begin{opts}
  \item 50
  \item 54
  \item 60
  \item 66
\end{opts}

\q{Con el impuesto de 5, la recaudación total es:}
\begin{opts}
  \item 270
  \item 300
  \item 240
  \item 150
\end{opts}

\q{Con el impuesto de 5, la pérdida de bienestar es:}
\begin{opts}
  \item 5
  \item 10
  \item 30
  \item 15
\end{opts}

\q{En el impuesto anterior, la carga que soportan los consumidores por unidad es:}
\begin{opts}
  \item 2
  \item 4
  \item 3
  \item 5
\end{opts}

\q{En el mismo impuesto, la carga que soportan los productores por unidad es:}
\begin{opts}
  \item 3
  \item 2
  \item 1
  \item 5
\end{opts}

\q{Si se aplica un subsidio de 5 por unidad, el nuevo precio de equilibrio es:}
\begin{opts}
  \item 27
  \item 30
  \item 33
  \item 25
\end{opts}

\q{Con el subsidio de 5, el costo total del subsidio es:}
\begin{opts}
  \item 330
  \item 270
  \item 300
  \item 360
\end{opts}

\q{Con el subsidio de 5, la pérdida de bienestar es:}
\begin{opts}
  \item 5
  \item 10
  \item 15
  \item 30
\end{opts}

\q{En una industria de costo constante, la oferta de largo plazo es:}
\begin{opts}
  \item ascendente.
  \item horizontal.
  \item descendente.
  \item vertical.
\end{opts}

\q{En una industria de costo creciente, la oferta de largo plazo es:}
\begin{opts}
  \item horizontal.
  \item descendente.
  \item ascendente.
  \item vertical.
\end{opts}

\section*{Bloque C. Casos con implicancias tarifarias y regulatorias del sector saneamiento}

\noindent\textit{Preguntas 66 a 90. Aplican los conceptos de competencia perfecta, equilibrio de mercado, excedentes, incidencia de impuestos y controles de precios a mercados de agua, servicios de saneamiento y políticas tarifarias.}

\q{En un mercado competitivo de camiones cisterna, la demanda es $Q_d=80-2P$ y la oferta es $Q_s=2P-20$. El precio de equilibrio es:}
\begin{opts}
    \item 20
    \item 25
    \item 30
    \item 35
\end{opts}

\q{En el mercado anterior, el excedente del consumidor es:}
\begin{opts}
    \item 112.5
    \item 225
    \item 450
    \item 900
\end{opts}

\q{En el mismo mercado, el excedente del productor es:}
\begin{opts}
    \item 112.5
    \item 450
    \item 900
    \item 225
\end{opts}

\q{En el mismo mercado, el excedente total es:}
\begin{opts}
    \item 225
    \item 900
    \item 450
    \item 675
\end{opts}

\q{Si el regulador fija un precio máximo de 20 en ese mercado, la escasez es:}
\begin{opts}
    \item 10
    \item 20
    \item 30
    \item 40
\end{opts}

\q{Con el precio máximo de 20, la pérdida de bienestar es:}
\begin{opts}
    \item 50
    \item 75
    \item 200
    \item 100
\end{opts}

\q{Si se aplica un impuesto específico de 4 por unidad, el nuevo precio de equilibrio es:}
\begin{opts}
    \item 25
    \item 27
    \item 29
    \item 31
\end{opts}

\q{Con el impuesto de 4, la recaudación total es:}
\begin{opts}
    \item 52
    \item 80
    \item 100
    \item 104
\end{opts}

\q{Con el impuesto de 4, la pérdida de bienestar es:}
\begin{opts}
    \item 8
    \item 16
    \item 4
    \item 12
\end{opts}

\q{En el impuesto anterior, la carga que soportan los consumidores por unidad es:}
\begin{opts}
    \item 1
    \item 3
    \item 4
    \item 2
\end{opts}

\q{En el mismo impuesto, la carga que soportan los productores por unidad es:}
\begin{opts}
    \item 1
    \item 2
    \item 3
    \item 4
\end{opts}

\q{Si se aplica un subsidio de 4 por unidad, el nuevo precio de equilibrio es:}
\begin{opts}
    \item 21
    \item 25
    \item 23
    \item 27
\end{opts}

\q{Con el subsidio de 4, el costo total del subsidio es:}
\begin{opts}
    \item 104
    \item 120
    \item 136
    \item 160
\end{opts}

\q{Con el subsidio de 4, la pérdida de bienestar es:}
\begin{opts}
    \item 8
    \item 16
    \item 4
    \item 12
\end{opts}

\q{Si en una industria competitiva el precio es mayor que el costo medio mínimo, en el largo plazo:}
\begin{opts}
    \item salen empresas de la industria y el precio tiende a subir.
    \item entran empresas, sube la oferta y el precio tiende a bajar.
    \item el equilibrio de mercado permanece completamente inalterado.
    \item el precio sigue subiendo de forma indefinida en el tiempo.
\end{opts}

\q{En una industria de costo constante, la oferta de largo plazo es:}
\begin{opts}
    \item horizontal.
    \item ascendente.
    \item descendente.
    \item vertical.
\end{opts}

\q{En un monopolio natural con costo medio decreciente, fijar el precio igual al costo marginal genera:}
\begin{opts}
    \item beneficios extraordinarios positivos.
    \item beneficio económico exactamente nulo.
    \item ausencia total de déficit financiero.
    \item pérdidas que exigen subsidio o cargo.
\end{opts}

\q{La teoría de mercados contestables sostiene que, incluso con pocas empresas, el precio puede aproximarse al competitivo si:}
\begin{opts}
    \item existen barreras legales infranqueables a la entrada.
    \item la amenaza de entrada de nuevos rivales es creíble.
    \item la empresa incumbente actúa como monopolio puro cerrado.
    \item las empresas establecidas fijan precios coordinadamente.
\end{opts}

\q{Si la tarifa de agua disminuye y todo lo demás permanece constante, el excedente del consumidor:}
\begin{opts}
    \item disminuye.
    \item aumenta.
    \item no cambia.
    \item se vuelve cero.
\end{opts}

\q{Si la tarifa de agua disminuye y todo lo demás permanece constante, el excedente del productor:}
\begin{opts}
    \item aumenta proporcionalmente con el volumen de agua.
    \item permanece exactamente igual en el corto plazo.
    \item disminuye y puede comprometer la sostenibilidad.
    \item se vuelve nulo aunque el precio siga sobre $CVMe$.
\end{opts}

\q{En un mercado de derechos de agua, la asignación competitiva tiende a:}
\begin{opts}
    \item repartir el agua por igual entre todos los usuarios.
    \item asignar el agua a los usos de mayor valor económico.
    \item asignar el agua a los usos de menor valor.
    \item distribuir el agua de forma aleatoria.
\end{opts}

\q{Un impuesto pigouviano sobre el consumo de agua que genera un costo externo debe fijarse igual a:}
\begin{opts}
    \item cero.
    \item el costo externo marginal.
    \item el costo medio total.
    \item el costo total del servicio.
\end{opts}

\q{Un precio máximo por debajo del equilibrio en un mercado de agua genera:}
\begin{opts}
    \item un excedente persistente de oferta.
    \item ningún efecto sobre las transacciones.
    \item escasez y racionamiento sin precios.
    \item un aumento automático de la oferta.
\end{opts}

\q{Si el objetivo es proteger a los hogares vulnerables sin generar escasez, la opción más eficiente es:}
\begin{opts}
    \item fijar un precio máximo general por debajo del equilibrio.
    \item establecer un precio mínimo obligatorio superior al equilibrio.
    \item imponer cuotas físicas uniformes de consumo por vivienda.
    \item aplicar un subsidio directo focalizado sin distorsionar precios.
\end{opts}

\q{El equilibrio competitivo se caracteriza por ser eficiente en sentido asignativo, pero:}
\begin{opts}
    \item garantiza siempre una distribución totalmente equitativa.
    \item resulta ineficiente incluso bajo información simétrica.
    \item prescinde de cualquier política pública redistributiva.
    \item puede requerir transferencias para lograr metas de equidad.
\end{opts}

\newpage
\section*{Tabla de claves rápidas}

\noindent Use esta tabla para una verificación inmediata. El solucionario con la justificación de cada respuesta se presenta en la sección siguiente.

\vspace{0.5cm}
\begin{center}
\renewcommand{\arraystretch}{1.35}
\begin{tabular}{|c|c|c|c|c|c|c|c|c|c|c|}
\hline
\textbf{N.} & \textbf{1} & \textbf{2} & \textbf{3} & \textbf{4} & \textbf{5} & \textbf{6} & \textbf{7} & \textbf{8} & \textbf{9} & \textbf{10} \\
\hline
\textbf{Clave} & C & A & C & D & C & D & B & D & B & D \\
\hline
\textbf{N.} & \textbf{11} & \textbf{12} & \textbf{13} & \textbf{14} & \textbf{15} & \textbf{16} & \textbf{17} & \textbf{18} & \textbf{19} & \textbf{20} \\
\hline
\textbf{Clave} & B & C & D & A & C & A & D & C & A & A \\
\hline
\textbf{N.} & \textbf{21} & \textbf{22} & \textbf{23} & \textbf{24} & \textbf{25} & \textbf{26} & \textbf{27} & \textbf{28} & \textbf{29} & \textbf{30} \\
\hline
\textbf{Clave} & C & A & A & B & C & C & D & D & A & A \\
\hline
\textbf{N.} & \textbf{31} & \textbf{32} & \textbf{33} & \textbf{34} & \textbf{35} & \textbf{36} & \textbf{37} & \textbf{38} & \textbf{39} & \textbf{40} \\
\hline
\textbf{Clave} & B & A & C & B & A & D & B & A & A & D \\
\hline
\textbf{N.} & \textbf{41} & \textbf{42} & \textbf{43} & \textbf{44} & \textbf{45} & \textbf{46} & \textbf{47} & \textbf{48} & \textbf{49} & \textbf{50} \\
\hline
\textbf{Clave} & B & C & A & C & B & C & A & B & C & D \\
\hline
\textbf{N.} & \textbf{51} & \textbf{52} & \textbf{53} & \textbf{54} & \textbf{55} & \textbf{56} & \textbf{57} & \textbf{58} & \textbf{59} & \textbf{60} \\
\hline
\textbf{Clave} & C & D & D & A & D & B & A & D & C & B \\
\hline
\textbf{N.} & \textbf{61} & \textbf{62} & \textbf{63} & \textbf{64} & \textbf{65} & \textbf{66} & \textbf{67} & \textbf{68} & \textbf{69} & \textbf{70} \\
\hline
\textbf{Clave} & A & A & C & B & C & B & B & D & C & B \\
\hline
\textbf{N.} & \textbf{71} & \textbf{72} & \textbf{73} & \textbf{74} & \textbf{75} & \textbf{76} & \textbf{77} & \textbf{78} & \textbf{79} & \textbf{80} \\
\hline
\textbf{Clave} & D & B & D & A & D & B & C & C & A & B \\
\hline
\textbf{N.} & \textbf{81} & \textbf{82} & \textbf{83} & \textbf{84} & \textbf{85} & \textbf{86} & \textbf{87} & \textbf{88} & \textbf{89} & \textbf{90} \\
\hline
\textbf{Clave} & A & D & B & B & C & B & B & C & D & D \\
\hline
\end{tabular}
\end{center}

\vspace{0.4cm}
\noindent\textbf{Distribución de claves:} A = 22; B = 23; C = 23; D = 22.

\newpage
\section*{Solucionario completo}

\noindent Cada respuesta incluye la clave y una justificación breve. En los ejercicios numéricos se indica la ecuación o el cálculo principal.

\subsection*{Bloque A. Fundamentos teóricos y conceptuales}

\sol{1}{C}{La competencia perfecta supone muchos compradores y vendedores, producto homogéneo, información perfecta, libre entrada y salida, y empresas precio-aceptantes.}

\sol{2}{A}{La empresa individual enfrenta una demanda horizontal: puede vender cualquier cantidad al precio de mercado, pero no puede fijar un precio superior.}

\sol{3}{C}{La demanda del mercado es descendente: a mayor precio, menor cantidad demandada por el conjunto de consumidores.}

\sol{4}{D}{La regla de maximización es $IMg=CMg$. En competencia perfecta, $IMg=P$, por lo que también $P=CMg$.}

\sol{5}{C}{Como la empresa vende cada unidad adicional al mismo precio, $IMg=P$.}

\sol{6}{D}{El beneficio económico es $IT-CT$ económico, donde $CT$ incluye costos explícitos e implícitos o de oportunidad.}

\sol{7}{B}{El beneficio normal remunera el costo de oportunidad del capital propio. Cuando el beneficio económico es cero, la empresa obtiene exactamente su beneficio normal.}

\sol{8}{D}{En el corto plazo, el beneficio puede ser positivo si $P>CMe$, cero si $P=CMe$, o negativo si $P<CMe$, siempre que decida producir.}

\sol{9}{B}{Si $P<CVMe$, la empresa no cubre ni los costos variables; le conviene suspender. Si $CVMe<P<CMe$, cubre variables y parte de los fijos.}

\sol{10}{D}{La oferta de corto plazo es el tramo de $CMg$ situado por encima del mínimo $CVMe$; por debajo de ese precio, la empresa no produce.}

\sol{11}{B}{La oferta de la industria es la suma horizontal de las cantidades ofrecidas por cada empresa a cada precio.}

\sol{12}{C}{Con libre entrada y salida, el equilibrio de largo plazo implica $P=CMg=CMe$ y beneficio económico nulo; las empresas solo obtienen beneficio normal.}

\sol{13}{D}{Si $P>\min CMe$, hay beneficios extraordinarios, entran empresas, la oferta aumenta y el precio cae hasta el mínimo del costo medio.}

\sol{14}{A}{Si $P<\min CMe$, las empresas incurren en pérdidas, salen del mercado, la oferta disminuye y el precio sube.}

\sol{15}{C}{En una industria de costo constante, los precios de los insumos no cambian al expandirse la industria; la oferta de largo plazo es horizontal.}

\sol{16}{A}{Si al expandirse la industria los costos de los insumos aumentan, la oferta de largo plazo tiene pendiente positiva.}

\sol{17}{D}{Si la expansión reduce los costos de los insumos, la oferta de largo plazo tiene pendiente negativa.}

\sol{18}{C}{El excedente del productor es $IT-CV$. También puede medirse como el área situada entre el precio y la curva de costo marginal.}

\sol{19}{A}{El excedente del consumidor es la diferencia entre la disposición a pagar y el precio efectivamente pagado.}

\sol{20}{A}{El excedente total es $EC+EP$ y se maximiza en el equilibrio competitivo, donde $P=CMg$.}

\sol{21}{C}{Un precio máximo por debajo del equilibrio reduce la cantidad ofrecida y aumenta la demandada: genera escasez y pérdida de bienestar.}

\sol{22}{A}{Un precio mínimo por encima del equilibrio genera excedente de oferta y pérdida de bienestar.}

\sol{23}{A}{La carga económica no depende solo de quién paga legalmente el impuesto, sino de las elasticidades: el lado menos elástico soporta una mayor proporción.}

\sol{24}{B}{En un subsidio, el beneficio se reparte según las elasticidades; el lado menos elástico captura una mayor parte.}

\sol{25}{C}{La eficiencia asignativa requiere que el precio refleje el costo marginal: $P=CMg$.}

\sol{26}{C}{En el largo plazo competitivo, la libre entrada lleva a producir en el mínimo del costo medio: $P=\min CMe$.}

\sol{27}{D}{La libre entrada elimina los beneficios extraordinarios: si hay beneficios, entran empresas hasta que $P=\min CMe$ y el beneficio económico se anula.}

\sol{28}{D}{Una empresa precio-aceptante no puede influir en el precio de mercado; lo toma como dato y decide solo su cantidad.}

\sol{29}{A}{El equilibrio de mercado se alcanza cuando la cantidad demandada iguala la cantidad ofrecida: $Q_d=Q_s$.}

\sol{30}{A}{En el corto plazo, la empresa produce donde $P=CMg$, siempre que el precio cubra al menos el costo variable medio.}

\subsection*{Bloque B. Aplicaciones cuantitativas y optimización}

\sol{31}{B}{$CMg=10+2q$. Igualando $P=30$: $10+2q=30$, de donde $q=10$.}

\sol{32}{A}{$IT=30(10)=300$. $CT=100+10(10)+10^2=300$. Beneficio $=300-300=0$.}

\sol{33}{C}{$CMg=5+q$. Igualando $P=25$: $5+q=25$, de donde $q=20$.}

\sol{34}{B}{$IT=25(20)=500$. $CT=50+5(20)+0.5(20)^2=50+100+200=350$. Beneficio $=150$.}

\sol{35}{A}{$CMg=4+2q=20$, entonces $q=8$. $IT=160$; $CT=100+4(8)+8^2=196$; beneficio $=-36$.}

\sol{36}{D}{$CVMe=4+q$, que se minimiza en $q=0$ con valor 4. Si $P<4$, la empresa no cubre los costos variables y suspende.}

\sol{37}{B}{$CMg=2q+10$. Igualando $P=30$: $2q+10=30$, de donde $q=10$.}

\sol{38}{A}{La oferta individual es $q=(P-10)/2$. Con 100 empresas: $Q=100(P-10)/2=50P-500$.}

\sol{39}{A}{$1000-50P=50P-500$, entonces $1500=100P$ y $P=15$.}

\sol{40}{D}{Sustituyendo $P=15$: $Q=1000-50(15)=250$.}

\sol{41}{B}{$q=(P-10)/2=(15-10)/2=2.5$. También $250/100=2.5$.}

\sol{42}{C}{$CMe=q+4+100/q$. Derivando: $1-100/q^2=0$, de donde $q=10$ y $CMe=10+4+10=24$. En largo plazo, $P=24$.}

\sol{43}{A}{El mínimo del costo medio se alcanza en $q=10$.}

\sol{44}{C}{Con $P=24$: $Q=500-10(24)=260$. Si cada empresa produce 10, el número de empresas es $260/10=26$.}

\sol{45}{B}{La demanda inversa es $P=50-0.1Q$. $EC=0.5(50-24)(260)=3380$.}

\sol{46}{C}{En el equilibrio de largo plazo con libre entrada, el beneficio económico de cada empresa es cero.}

\sol{47}{A}{En una industria de costo constante, la oferta de largo plazo es horizontal al mínimo del costo medio. Un aumento de demanda eleva la cantidad, pero el precio vuelve a 24.}

\sol{48}{B}{$120-2P=3P-30$, entonces $150=5P$ y $P=30$.}

\sol{49}{C}{Sustituyendo $P=30$: $Q=120-2(30)=60$.}

\sol{50}{D}{La demanda inversa es $P=60-0.5Q$. $EC=0.5(60-30)(60)=900$.}

\sol{51}{C}{La oferta inversa es $P=10+Q/3$. $EP=0.5(30-10)(60)=600$.}

\sol{52}{D}{$ET=EC+EP=900+600=1500$.}

\sol{53}{D}{Con $P=20$: $Q_d=120-40=80$ y $Q_s=60-30=30$. La escasez es $80-30=50$.}

\sol{54}{A}{En $Q=30$, el precio de demanda es $(120-30)/2=45$. La pérdida es $0.5(80-30)(45-20)=625$.}

\sol{55}{D}{Con impuesto de 5, la oferta es $Q_s=3(P-5)-30=3P-45$. Igualando: $120-2P=3P-45$, de donde $P=33$.}

\sol{56}{B}{Con $P=33$: $Q=120-2(33)=54$.}

\sol{57}{A}{Recaudación $=5(54)=270$.}

\sol{58}{D}{$PCP=0.5(5)(60-54)=15$.}

\sol{59}{C}{El precio sube de 30 a 33: los consumidores soportan 3 por unidad.}

\sol{60}{B}{El productor recibe $33-5=28$, frente a 30 inicial: soporta 2 por unidad.}

\sol{61}{A}{Con subsidio de 5, la oferta es $Q_s=3(P+5)-30=3P-15$. Igualando: $120-2P=3P-15$, de donde $P=27$.}

\sol{62}{A}{Con $P=27$: $Q=120-2(27)=66$. Costo del subsidio $=5(66)=330$.}

\sol{63}{C}{$PCP=0.5(5)(66-60)=15$.}

\sol{64}{B}{En una industria de costo constante, la oferta de largo plazo es horizontal.}

\sol{65}{C}{En una industria de costo creciente, la oferta de largo plazo es ascendente.}

\subsection*{Bloque C. Casos con implicancias tarifarias y regulatorias}

\sol{66}{B}{$80-2P=2P-20$, entonces $100=4P$ y $P=25$.}

\sol{67}{B}{La demanda inversa es $P=40-0.5Q$. $EC=0.5(40-25)(30)=225$.}

\sol{68}{D}{La oferta inversa es $P=10+0.5Q$. $EP=0.5(25-10)(30)=225$.}

\sol{69}{C}{$ET=225+225=450$.}

\sol{70}{B}{Con $P=20$: $Q_d=80-40=40$ y $Q_s=40-20=20$. La escasez es 20.}

\sol{71}{D}{En $Q=20$, el precio de demanda es $(80-20)/2=30$. $PCP=0.5(40-20)(30-20)=100$.}

\sol{72}{B}{Con impuesto de 4, la oferta es $Q_s=2(P-4)-20=2P-28$. Igualando: $80-2P=2P-28$, de donde $P=27$.}

\sol{73}{D}{Con $P=27$: $Q=80-2(27)=26$. Recaudación $=4(26)=104$.}

\sol{74}{A}{$PCP=0.5(4)(30-26)=8$.}

\sol{75}{D}{El precio sube de 25 a 27: los consumidores soportan 2 por unidad.}

\sol{76}{B}{El productor recibe $27-4=23$, frente a 25 inicial: soporta 2 por unidad.}

\sol{77}{C}{Con subsidio de 4, la oferta es $Q_s=2(P+4)-20=2P-12$. Igualando: $80-2P=2P-12$, de donde $P=23$.}

\sol{78}{C}{Con $P=23$: $Q=80-2(23)=34$. Costo del subsidio $=4(34)=136$.}

\sol{79}{A}{$PCP=0.5(4)(34-30)=8$.}

\sol{80}{B}{Si $P>\min CMe$, hay beneficios extraordinarios; entran empresas, la oferta aumenta y el precio tiende a bajar hasta el mínimo del costo medio.}

\sol{81}{A}{En una industria de costo constante, la oferta de largo plazo es horizontal.}

\sol{82}{D}{En un monopolio natural, $P=CMg<CMe$ genera pérdidas; se requiere un subsidio o un cargo fijo para cubrirlas.}

\sol{83}{B}{Un mercado contestable puede aproximarse al resultado competitivo si la entrada potencial es creíble, aunque haya pocas empresas.}

\sol{84}{B}{Una reducción de la tarifa aumenta el excedente del consumidor, pues paga menos por la misma cantidad o consume más.}

\sol{85}{C}{Una reducción de la tarifa disminuye el excedente del productor y puede comprometer la recuperación de costos y la sostenibilidad financiera.}

\sol{86}{B}{Un mercado de derechos de agua permite que el recurso se reasigne hacia los usos de mayor valor económico, sujeto a los límites ambientales.}

\sol{87}{B}{El impuesto pigouviano óptimo iguala el costo externo marginal, internalizando la externalidad negativa.}

\sol{88}{C}{Un precio máximo por debajo del equilibrio genera escasez; la asignación se realiza por mecanismos no basados en el precio, como racionamiento o colas.}

\sol{89}{D}{Un subsidio directo focalizado protege el ingreso de los hogares vulnerables sin distorsionar el precio de mercado ni generar escasez.}

\sol{90}{D}{El equilibrio competitivo es eficiente en sentido asignativo, pero no garantiza equidad; pueden requerirse transferencias o subsidios focalizados.}


\end{document}
